% Options for packages loaded elsewhere
\PassOptionsToPackage{unicode}{hyperref}
\PassOptionsToPackage{hyphens}{url}
\documentclass[
  spanish,
]{article}
\usepackage{xcolor}
\usepackage[margin=1in]{geometry}
\usepackage{amsmath,amssymb}
\setcounter{secnumdepth}{5}
\usepackage{iftex}
\ifPDFTeX
  \usepackage[T1]{fontenc}
  \usepackage[utf8]{inputenc}
  \usepackage{textcomp} % provide euro and other symbols
\else % if luatex or xetex
  \usepackage{unicode-math} % this also loads fontspec
  \defaultfontfeatures{Scale=MatchLowercase}
  \defaultfontfeatures[\rmfamily]{Ligatures=TeX,Scale=1}
\fi
\usepackage{lmodern}
\ifPDFTeX\else
  % xetex/luatex font selection
\fi
% Use upquote if available, for straight quotes in verbatim environments
\IfFileExists{upquote.sty}{\usepackage{upquote}}{}
\IfFileExists{microtype.sty}{% use microtype if available
  \usepackage[]{microtype}
  \UseMicrotypeSet[protrusion]{basicmath} % disable protrusion for tt fonts
}{}
\makeatletter
\@ifundefined{KOMAClassName}{% if non-KOMA class
  \IfFileExists{parskip.sty}{%
    \usepackage{parskip}
  }{% else
    \setlength{\parindent}{0pt}
    \setlength{\parskip}{6pt plus 2pt minus 1pt}}
}{% if KOMA class
  \KOMAoptions{parskip=half}}
\makeatother
\usepackage{graphicx}
\makeatletter
\newsavebox\pandoc@box
\newcommand*\pandocbounded[1]{% scales image to fit in text height/width
  \sbox\pandoc@box{#1}%
  \Gscale@div\@tempa{\textheight}{\dimexpr\ht\pandoc@box+\dp\pandoc@box\relax}%
  \Gscale@div\@tempb{\linewidth}{\wd\pandoc@box}%
  \ifdim\@tempb\p@<\@tempa\p@\let\@tempa\@tempb\fi% select the smaller of both
  \ifdim\@tempa\p@<\p@\scalebox{\@tempa}{\usebox\pandoc@box}%
  \else\usebox{\pandoc@box}%
  \fi%
}
% Set default figure placement to htbp
\def\fps@figure{htbp}
\makeatother
\ifLuaTeX
\usepackage[bidi=basic]{babel}
\else
\usepackage[bidi=default]{babel}
\fi
% get rid of language-specific shorthands (see #6817):
\let\LanguageShortHands\languageshorthands
\def\languageshorthands#1{}
\setlength{\emergencystretch}{3em} % prevent overfull lines
\providecommand{\tightlist}{%
  \setlength{\itemsep}{0pt}\setlength{\parskip}{0pt}}
\usepackage{booktabs}
\usepackage{longtable}
\usepackage{array}
\usepackage{float}
\usepackage{colortbl}
\usepackage{xcolor}
\usepackage{geometry}
\geometry{margin=2.5cm}
\usepackage{fancyhdr}
\pagestyle{fancy}
Factor Sensorial}
\fancyhead[R]{\small Tesis --- Experiencia Gastronómica}
\fancyfoot[C]{\thepage}
\renewcommand{\headrulewidth}{0.4pt}
\usepackage{caption}
\captionsetup{font=small, labelfont=bf}
\usepackage{indentfirst}
\setlength{\parindent}{1.25cm}
\usepackage{booktabs}
\usepackage{longtable}
\usepackage{array}
\usepackage{multirow}
\usepackage{wrapfig}
\usepackage{float}
\usepackage{colortbl}
\usepackage{pdflscape}
\usepackage{tabu}
\usepackage{threeparttable}
\usepackage{threeparttablex}
\usepackage[normalem]{ulem}
\usepackage{makecell}
\usepackage{xcolor}
\usepackage{bookmark}
\IfFileExists{xurl.sty}{\usepackage{xurl}}{} % add URL line breaks if available
\urlstyle{same}
\hypersetup{
  pdftitle={Objetivo 4: Factor Sensorial Más Influyente},
  pdfauthor={Ariana-Fernanda Valeria Vidal Tapia},
  pdflang={es},
  hidelinks,
  pdfcreator={LaTeX via pandoc}}

\title{Objetivo 4: Factor Sensorial Más Influyente}
\usepackage{etoolbox}
\makeatletter
\providecommand{\subtitle}[1]{% add subtitle to \maketitle
  \apptocmd{\@title}{\par {\large #1 \par}}{}{}
}
\makeatother
\subtitle{Impacto de los factores de la experiencia gastronómica en la
toma de decisiones del comensal}
\author{Ariana-Fernanda Valeria Vidal Tapia}
\date{16 de April de 2026}

\begin{document}
\maketitle

{
\setcounter{tocdepth}{3}
\tableofcontents
}
\section{Carga de datos y
librerías}\label{carga-de-datos-y-libreruxedas}

El objetivo específico 4 busca identificar cuál de los factores
sensoriales dentro de cada dimensión --- Producto (P1, P2, P3), Servicio
(S1, S2, S3) y Entorno (E1--E8) --- es el más influyente en la toma de
decisiones del comensal, medida a través de tres variables dependientes:
satisfacción (I1, Likert 1--5), intención de revisitar (I2, Likert 1--5)
e intención de recomendar (I3, binaria Sí/No).

Para cada dimensión se presentan: correlaciones de Spearman con ranking
de influencia, tablas de contingencia, y visualizaciones (heatmap,
barras apiladas y dispersión con jitter).

\section{Tablas de puntajes promedio por
condición}\label{tablas-de-puntajes-promedio-por-condiciuxf3n}

\begin{longtable}[t]{lccc}
\caption{\label{tab:tabla-calidad-producto}Calidad de Producto}\\
\toprule
Condición & P1 (Sabor) & P2 (Presentación) & P3 (Frescura)\\
\midrule
Base (todo bien) & 3.61 & 3.89 & 3.80\\
Servicio degradado & 3.09 & 3.70 & 3.52\\
Producto degradado & 2.36 & 3.07 & 2.73\\
Entorno degradado & 3.11 & 3.02 & 3.32\\
\bottomrule
\end{longtable}

\begin{longtable}[t]{lccc}
\caption{\label{tab:tabla-calidad-servicio}Calidad del Servicio}\\
\toprule
Condición & S1 (Atención) & S2 (Confianza) & S3 (Apariencia)\\
\midrule
Base (todo bien) & 4.39 & 4.16 & 4.36\\
Servicio degradado & 2.80 & 2.91 & 3.61\\
Producto degradado & 4.30 & 4.34 & 4.34\\
Entorno degradado & 3.82 & 3.82 & 4.20\\
\bottomrule
\end{longtable}

\begingroup\fontsize{9}{11}\selectfont

\begin{longtable}[t]{lcccccccc}
\caption{\label{tab:tabla-calidad-entorno}Calidad del Entorno}\\
\toprule
Condición & E1 (Decoración) & E2 (Colores) & E3 (Mobiliario) & E4 (Menaje) & E5 (Manteles) & E6 (Espacio) & E7 (Asientos) & E8 (Música)\\
\midrule
Base (todo bien) & 4.70 & 4.45 & 3.95 & 3.98 & 4.30 & 4.52 & 4.25 & 4.45\\
Servicio degradado & 4.36 & 4.32 & 3.66 & 3.86 & 4.16 & 4.36 & 4.36 & 3.75\\
Producto degradado & 4.41 & 4.39 & 3.73 & 3.89 & 4.16 & 4.34 & 4.39 & 4.02\\
Entorno degradado & 2.30 & 2.32 & 2.64 & 3.11 & 3.30 & 1.82 & 1.59 & 1.27\\
\bottomrule
\end{longtable}
\endgroup{}

\newpage

\section{Análisis por ítems de Producto: P1, P2, P3 vs.~variables
dependientes}\label{anuxe1lisis-por-uxedtems-de-producto-p1-p2-p3-vs.-variables-dependientes}

\subsection{Correlaciones de Spearman}\label{correlaciones-de-spearman}

Se calcula el coeficiente de correlación de Spearman (\(\rho_s\)) entre
cada ítem del producto y cada variable dependiente.

\begin{longtable}[t]{llcc>{}c}
\toprule
Ítem Producto & Variable dependiente & \$\textbackslash{}rho\_s\$ & \$p\$ & Significativo\\
\midrule
\addlinespace[0.3em]
\multicolumn{5}{l}{\textbf{P1 (Sabor)}}\\
\hspace{1em}P1 (Sabor) & Satisfacción & 0.586 & < .001 & \textcolor{teal}{Sí}\\
\hspace{1em}P1 (Sabor) & Revisitar & 0.550 & < .001 & \textcolor{teal}{Sí}\\
\hspace{1em}P1 (Sabor) & Recomendar & 0.484 & < .001 & \textcolor{teal}{Sí}\\
\addlinespace[0.3em]
\multicolumn{5}{l}{\textbf{P2 (Presentación)}}\\
\hspace{1em}P2 (Presentación) & Satisfacción & 0.490 & < .001 & \textcolor{teal}{Sí}\\
\hspace{1em}P2 (Presentación) & Revisitar & 0.517 & < .001 & \textcolor{teal}{Sí}\\
\hspace{1em}P2 (Presentación) & Recomendar & 0.427 & < .001 & \textcolor{teal}{Sí}\\
\addlinespace[0.3em]
\multicolumn{5}{l}{\textbf{P3 (Frescura)}}\\
\hspace{1em}P3 (Frescura) & Satisfacción & 0.414 & < .001 & \textcolor{teal}{Sí}\\
\hspace{1em}P3 (Frescura) & Revisitar & 0.351 & < .001 & \textcolor{teal}{Sí}\\
\hspace{1em}P3 (Frescura) & Recomendar & 0.364 & < .001 & \textcolor{teal}{Sí}\\
\bottomrule
\end{longtable}

\subsubsection{Ranking de influencia
(Producto)}\label{ranking-de-influencia-producto}

\begin{longtable}[t]{clc}
\toprule
Ranking & Ítem Producto & \$\&\#124;\textbackslash{}rho\_s\&\#124;\$ promedio\\
\midrule
1 & P1 (Sabor) & 0.540\\
2 & P2 (Presentación) & 0.478\\
3 & P3 (Frescura) & 0.376\\
\bottomrule
\end{longtable}

\newpage

\subsection{Tablas de contingencia
(Producto)}\label{tablas-de-contingencia-producto}

\subsubsection{P1 (Sabor) vs.~Variables
dependientes}\label{p1-sabor-vs.-variables-dependientes}

\begin{longtable}[t]{lrrrrrr}
\caption{\label{tab:cont-p1-satisf}P1 (Sabor) vs. Satisfacción (I1)}\\
\toprule
 & 1 & 2 & 3 & 4 & 5 & Sum\\
\midrule
1 & 3 & 6 & 1 & 0 & 0 & 10\\
2 & 2 & 21 & 12 & 4 & 0 & 39\\
3 & 1 & 7 & 25 & 37 & 0 & 70\\
4 & 0 & 3 & 7 & 29 & 8 & 47\\
5 & 0 & 2 & 2 & 4 & 2 & 10\\
\addlinespace
Sum & 6 & 39 & 47 & 74 & 10 & 176\\
\bottomrule
\end{longtable}

\begin{longtable}[t]{lrrrrrr}
\caption{\label{tab:cont-p1-revis}P1 (Sabor) vs. Revisitar (I2)}\\
\toprule
 & 1 & 2 & 3 & 4 & 5 & Sum\\
\midrule
1 & 5 & 4 & 1 & 0 & 0 & 10\\
2 & 2 & 23 & 13 & 1 & 0 & 39\\
3 & 1 & 10 & 34 & 23 & 2 & 70\\
4 & 0 & 7 & 8 & 22 & 10 & 47\\
5 & 1 & 1 & 4 & 1 & 3 & 10\\
\addlinespace
Sum & 9 & 45 & 60 & 47 & 15 & 176\\
\bottomrule
\end{longtable}

\begin{longtable}[t]{lrrr}
\caption{\label{tab:cont-p1-recom}P1 (Sabor) vs. Recomendar (I3)}\\
\toprule
 & No recomendaría & Sí recomendaría & Sum\\
\midrule
1 & 10 & 0 & 10\\
2 & 33 & 6 & 39\\
3 & 30 & 40 & 70\\
4 & 12 & 35 & 47\\
5 & 2 & 8 & 10\\
\addlinespace
Sum & 87 & 89 & 176\\
\bottomrule
\end{longtable}

\subsubsection{P2 (Presentación) vs.~Variables
dependientes}\label{p2-presentaciuxf3n-vs.-variables-dependientes}

\begin{longtable}[t]{lrrrrrr}
\caption{\label{tab:cont-p2-satisf}P2 (Presentación) vs. Satisfacción (I1)}\\
\toprule
 & 1 & 2 & 3 & 4 & 5 & Sum\\
\midrule
1 & 0 & 4 & 0 & 0 & 0 & 4\\
2 & 4 & 8 & 9 & 1 & 0 & 22\\
3 & 1 & 18 & 24 & 25 & 1 & 69\\
4 & 1 & 6 & 11 & 36 & 4 & 58\\
5 & 0 & 3 & 3 & 12 & 5 & 23\\
\addlinespace
Sum & 6 & 39 & 47 & 74 & 10 & 176\\
\bottomrule
\end{longtable}

\begin{longtable}[t]{lrrrrrr}
\caption{\label{tab:cont-p2-revis}P2 (Presentación) vs. Revisitar (I2)}\\
\toprule
 & 1 & 2 & 3 & 4 & 5 & Sum\\
\midrule
1 & 2 & 2 & 0 & 0 & 0 & 4\\
2 & 4 & 13 & 5 & 0 & 0 & 22\\
3 & 1 & 23 & 26 & 18 & 1 & 69\\
4 & 1 & 6 & 22 & 21 & 8 & 58\\
5 & 1 & 1 & 7 & 8 & 6 & 23\\
\addlinespace
Sum & 9 & 45 & 60 & 47 & 15 & 176\\
\bottomrule
\end{longtable}

\begin{longtable}[t]{lrrr}
\caption{\label{tab:cont-p2-recom}P2 (Presentación) vs. Recomendar (I3)}\\
\toprule
 & No recomendaría & Sí recomendaría & Sum\\
\midrule
1 & 4 & 0 & 4\\
2 & 20 & 2 & 22\\
3 & 39 & 30 & 69\\
4 & 19 & 39 & 58\\
5 & 5 & 18 & 23\\
\addlinespace
Sum & 87 & 89 & 176\\
\bottomrule
\end{longtable}

\subsubsection{P3 (Frescura) vs.~Variables
dependientes}\label{p3-frescura-vs.-variables-dependientes}

\begin{longtable}[t]{lrrrrrr}
\caption{\label{tab:cont-p3-satisf}P3 (Frescura) vs. Satisfacción (I1)}\\
\toprule
 & 1 & 2 & 3 & 4 & 5 & Sum\\
\midrule
1 & 2 & 6 & 0 & 0 & 0 & 8\\
2 & 1 & 9 & 7 & 10 & 0 & 27\\
3 & 2 & 12 & 29 & 18 & 2 & 63\\
4 & 1 & 10 & 8 & 31 & 3 & 53\\
5 & 0 & 2 & 3 & 15 & 5 & 25\\
\addlinespace
Sum & 6 & 39 & 47 & 74 & 10 & 176\\
\bottomrule
\end{longtable}

\begin{longtable}[t]{lrrrrrr}
\caption{\label{tab:cont-p3-revis}P3 (Frescura) vs. Revisitar (I2)}\\
\toprule
 & 1 & 2 & 3 & 4 & 5 & Sum\\
\midrule
1 & 4 & 2 & 2 & 0 & 0 & 8\\
2 & 1 & 11 & 10 & 4 & 1 & 27\\
3 & 1 & 19 & 24 & 16 & 3 & 63\\
4 & 2 & 12 & 17 & 17 & 5 & 53\\
5 & 1 & 1 & 7 & 10 & 6 & 25\\
\addlinespace
Sum & 9 & 45 & 60 & 47 & 15 & 176\\
\bottomrule
\end{longtable}

\begin{longtable}[t]{lrrr}
\caption{\label{tab:cont-p3-recom}P3 (Frescura) vs. Recomendar (I3)}\\
\toprule
 & No recomendaría & Sí recomendaría & Sum\\
\midrule
1 & 8 & 0 & 8\\
2 & 19 & 8 & 27\\
3 & 34 & 29 & 63\\
4 & 22 & 31 & 53\\
5 & 4 & 21 & 25\\
\addlinespace
Sum & 87 & 89 & 176\\
\bottomrule
\end{longtable}

\newpage

\subsection{Visualizaciones (Producto)}\label{visualizaciones-producto}

\subsubsection{Heatmap de correlaciones}\label{heatmap-de-correlaciones}

\begin{figure}[H]

{\centering \includegraphics[width=0.95\linewidth]{03_objetivo4_factor_sensorial-06-_files/figure-latex/heatmap-corr-1} 

}

\caption{Correlaciones de Spearman entre los ítems del Producto y las variables dependientes.}\label{fig:heatmap-corr}
\end{figure}

\subsubsection{Proporción de Recomendar por nivel de cada
ítem}\label{proporciuxf3n-de-recomendar-por-nivel-de-cada-uxedtem}

\begin{figure}[H]

{\centering \includegraphics[width=0.95\linewidth]{03_objetivo4_factor_sensorial-06-_files/figure-latex/mosaico-p1-1} 

}

\caption{Proporción de Recomendar por nivel de P1 (Sabor).}\label{fig:mosaico-p1}
\end{figure}

\begin{figure}[H]

{\centering \includegraphics[width=0.95\linewidth]{03_objetivo4_factor_sensorial-06-_files/figure-latex/mosaico-p2-1} 

}

\caption{Proporción de Recomendar por nivel de P2 (Presentación).}\label{fig:mosaico-p2}
\end{figure}

\begin{figure}[H]

{\centering \includegraphics[width=0.95\linewidth]{03_objetivo4_factor_sensorial-06-_files/figure-latex/mosaico-p3-1} 

}

\caption{Proporción de Recomendar por nivel de P3 (Frescura).}\label{fig:mosaico-p3}
\end{figure}

\subsubsection{Dispersión con jitter: Ítems del producto
vs.~Satisfacción}\label{dispersiuxf3n-con-jitter-uxedtems-del-producto-vs.-satisfacciuxf3n}

\begin{figure}[H]

{\centering \includegraphics[width=0.95\linewidth]{03_objetivo4_factor_sensorial-06-_files/figure-latex/jitter-satisf-1} 

}

\caption{Relación entre cada ítem del Producto y la Satisfacción (I1).}\label{fig:jitter-satisf}
\end{figure}

\subsubsection{Dispersión con jitter: Ítems del producto
vs.~Revisitar}\label{dispersiuxf3n-con-jitter-uxedtems-del-producto-vs.-revisitar}

\begin{figure}[H]

{\centering \includegraphics[width=0.95\linewidth]{03_objetivo4_factor_sensorial-06-_files/figure-latex/jitter-revisitar-1} 

}

\caption{Relación entre cada ítem del Producto y la Intención de revisitar (I2).}\label{fig:jitter-revisitar}
\end{figure}

\subsubsection{Resumen (Producto)}\label{resumen-producto}

El ítem de producto más influyente en la toma de decisiones del comensal
es \textbf{P1 (Sabor)}, con \(|\rho_s|\) promedio = 0.54.

\newpage

\section{Análisis por ítems de Servicio: S1, S2, S3 vs.~variables
dependientes}\label{anuxe1lisis-por-uxedtems-de-servicio-s1-s2-s3-vs.-variables-dependientes}

\subsection{Correlaciones de
Spearman}\label{correlaciones-de-spearman-1}

\begin{longtable}[t]{llcc>{}c}
\toprule
Ítem Servicio & Variable dependiente & \$\textbackslash{}rho\_s\$ & \$p\$ & Significativo\\
\midrule
\addlinespace[0.3em]
\multicolumn{5}{l}{\textbf{S1 (Atención)}}\\
\hspace{1em}S1 (Atención) & Satisfacción & 0.341 & < .001 & \textcolor{teal}{Sí}\\
\hspace{1em}S1 (Atención) & Revisitar & 0.273 & < .001 & \textcolor{teal}{Sí}\\
\hspace{1em}S1 (Atención) & Recomendar & 0.286 & < .001 & \textcolor{teal}{Sí}\\
\addlinespace[0.3em]
\multicolumn{5}{l}{\textbf{S2 (Confianza)}}\\
\hspace{1em}S2 (Confianza) & Satisfacción & 0.250 & < .001 & \textcolor{teal}{Sí}\\
\hspace{1em}S2 (Confianza) & Revisitar & 0.199 & 0.0081 & \textcolor{teal}{Sí}\\
\hspace{1em}S2 (Confianza) & Recomendar & 0.235 & 0.0017 & \textcolor{teal}{Sí}\\
\addlinespace[0.3em]
\multicolumn{5}{l}{\textbf{S3 (Apariencia)}}\\
\hspace{1em}S3 (Apariencia) & Satisfacción & 0.300 & < .001 & \textcolor{teal}{Sí}\\
\hspace{1em}S3 (Apariencia) & Revisitar & 0.198 & 0.0086 & \textcolor{teal}{Sí}\\
\hspace{1em}S3 (Apariencia) & Recomendar & 0.118 & 0.1176 & \textcolor{red}{No}\\
\bottomrule
\end{longtable}

\subsubsection{Ranking de influencia
(Servicio)}\label{ranking-de-influencia-servicio}

\begin{longtable}[t]{clc}
\toprule
Ranking & Ítem Servicio & \$\&\#124;\textbackslash{}rho\_s\&\#124;\$ promedio\\
\midrule
1 & S1 (Atención) & 0.300\\
2 & S2 (Confianza) & 0.228\\
3 & S3 (Apariencia) & 0.205\\
\bottomrule
\end{longtable}

\newpage

\subsection{Tablas de contingencia
(Servicio)}\label{tablas-de-contingencia-servicio}

\subsubsection{S1 (Atención) vs.~Variables
dependientes}\label{s1-atenciuxf3n-vs.-variables-dependientes}

\begin{longtable}[t]{lrrrrrr}
\caption{\label{tab:cont-s1-satisf}S1 (Atención) vs. Satisfacción (I1)}\\
\toprule
 & 1 & 2 & 3 & 4 & 5 & Sum\\
\midrule
1 & 1 & 2 & 2 & 0 & 0 & 5\\
2 & 0 & 7 & 9 & 6 & 0 & 22\\
3 & 1 & 11 & 7 & 12 & 0 & 31\\
4 & 2 & 15 & 16 & 23 & 3 & 59\\
5 & 2 & 4 & 13 & 33 & 7 & 59\\
\addlinespace
Sum & 6 & 39 & 47 & 74 & 10 & 176\\
\bottomrule
\end{longtable}

\begin{longtable}[t]{lrrrrrr}
\caption{\label{tab:cont-s1-revis}S1 (Atención) vs. Revisitar (I2)}\\
\toprule
 & 1 & 2 & 3 & 4 & 5 & Sum\\
\midrule
1 & 2 & 1 & 1 & 1 & 0 & 5\\
2 & 0 & 9 & 9 & 4 & 0 & 22\\
3 & 2 & 10 & 9 & 9 & 1 & 31\\
4 & 4 & 17 & 19 & 15 & 4 & 59\\
5 & 1 & 8 & 22 & 18 & 10 & 59\\
\addlinespace
Sum & 9 & 45 & 60 & 47 & 15 & 176\\
\bottomrule
\end{longtable}

\begin{longtable}[t]{lrrr}
\caption{\label{tab:cont-s1-recom}S1 (Atención) vs. Recomendar (I3)}\\
\toprule
 & No recomendaría & Sí recomendaría & Sum\\
\midrule
1 & 4 & 1 & 5\\
2 & 14 & 8 & 22\\
3 & 20 & 11 & 31\\
4 & 31 & 28 & 59\\
5 & 18 & 41 & 59\\
\addlinespace
Sum & 87 & 89 & 176\\
\bottomrule
\end{longtable}

\subsubsection{S2 (Confianza) vs.~Variables
dependientes}\label{s2-confianza-vs.-variables-dependientes}

\begin{longtable}[t]{lrrrrrr}
\caption{\label{tab:cont-s2-satisf}S2 (Confianza) vs. Satisfacción (I1)}\\
\toprule
 & 1 & 2 & 3 & 4 & 5 & Sum\\
\midrule
1 & 1 & 2 & 1 & 1 & 0 & 5\\
2 & 0 & 8 & 6 & 4 & 0 & 18\\
3 & 1 & 9 & 13 & 18 & 0 & 41\\
4 & 3 & 11 & 14 & 22 & 4 & 54\\
5 & 1 & 9 & 13 & 29 & 6 & 58\\
\addlinespace
Sum & 6 & 39 & 47 & 74 & 10 & 176\\
\bottomrule
\end{longtable}

\begin{longtable}[t]{lrrrrrr}
\caption{\label{tab:cont-s2-revis}S2 (Confianza) vs. Revisitar (I2)}\\
\toprule
 & 1 & 2 & 3 & 4 & 5 & Sum\\
\midrule
1 & 2 & 2 & 0 & 1 & 0 & 5\\
2 & 0 & 8 & 7 & 3 & 0 & 18\\
3 & 2 & 11 & 14 & 11 & 3 & 41\\
4 & 4 & 13 & 16 & 18 & 3 & 54\\
5 & 1 & 11 & 23 & 14 & 9 & 58\\
\addlinespace
Sum & 9 & 45 & 60 & 47 & 15 & 176\\
\bottomrule
\end{longtable}

\begin{longtable}[t]{lrrr}
\caption{\label{tab:cont-s2-recom}S2 (Confianza) vs. Recomendar (I3)}\\
\toprule
 & No recomendaría & Sí recomendaría & Sum\\
\midrule
1 & 4 & 1 & 5\\
2 & 11 & 7 & 18\\
3 & 25 & 16 & 41\\
4 & 27 & 27 & 54\\
5 & 20 & 38 & 58\\
\addlinespace
Sum & 87 & 89 & 176\\
\bottomrule
\end{longtable}

\subsubsection{S3 (Apariencia) vs.~Variables
dependientes}\label{s3-apariencia-vs.-variables-dependientes}

\begin{longtable}[t]{lrrrrrr}
\caption{\label{tab:cont-s3-satisf}S3 (Apariencia) vs. Satisfacción (I1)}\\
\toprule
 & 1 & 2 & 3 & 4 & 5 & Sum\\
\midrule
1 & 1 & 0 & 1 & 0 & 0 & 2\\
2 & 1 & 1 & 4 & 2 & 0 & 8\\
3 & 1 & 11 & 8 & 6 & 1 & 27\\
4 & 2 & 15 & 21 & 27 & 2 & 67\\
5 & 1 & 12 & 13 & 39 & 7 & 72\\
\addlinespace
Sum & 6 & 39 & 47 & 74 & 10 & 176\\
\bottomrule
\end{longtable}

\begin{longtable}[t]{lrrrrrr}
\caption{\label{tab:cont-s3-revis}S3 (Apariencia) vs. Revisitar (I2)}\\
\toprule
 & 1 & 2 & 3 & 4 & 5 & Sum\\
\midrule
1 & 1 & 0 & 0 & 1 & 0 & 2\\
2 & 2 & 1 & 3 & 2 & 0 & 8\\
3 & 1 & 12 & 9 & 5 & 0 & 27\\
4 & 3 & 17 & 22 & 20 & 5 & 67\\
5 & 2 & 15 & 26 & 19 & 10 & 72\\
\addlinespace
Sum & 9 & 45 & 60 & 47 & 15 & 176\\
\bottomrule
\end{longtable}

\begin{longtable}[t]{lrrr}
\caption{\label{tab:cont-s3-recom}S3 (Apariencia) vs. Recomendar (I3)}\\
\toprule
 & No recomendaría & Sí recomendaría & Sum\\
\midrule
1 & 1 & 1 & 2\\
2 & 4 & 4 & 8\\
3 & 17 & 10 & 27\\
4 & 34 & 33 & 67\\
5 & 31 & 41 & 72\\
\addlinespace
Sum & 87 & 89 & 176\\
\bottomrule
\end{longtable}

\newpage

\subsection{Visualizaciones (Servicio)}\label{visualizaciones-servicio}

\subsubsection{Heatmap de
correlaciones}\label{heatmap-de-correlaciones-1}

\begin{figure}[H]

{\centering \includegraphics[width=0.95\linewidth]{03_objetivo4_factor_sensorial-06-_files/figure-latex/heatmap-serv-1} 

}

\caption{Correlaciones de Spearman entre los ítems del Servicio y las variables dependientes.}\label{fig:heatmap-serv}
\end{figure}

\subsubsection{Proporción de Recomendar por nivel de cada
ítem}\label{proporciuxf3n-de-recomendar-por-nivel-de-cada-uxedtem-1}

\begin{figure}[H]

{\centering \includegraphics[width=0.95\linewidth]{03_objetivo4_factor_sensorial-06-_files/figure-latex/mosaico-s1-1} 

}

\caption{Proporción de Recomendar por nivel de S1 (Atención).}\label{fig:mosaico-s1}
\end{figure}

\begin{figure}[H]

{\centering \includegraphics[width=0.95\linewidth]{03_objetivo4_factor_sensorial-06-_files/figure-latex/mosaico-s2-1} 

}

\caption{Proporción de Recomendar por nivel de S2 (Confianza).}\label{fig:mosaico-s2}
\end{figure}

\begin{figure}[H]

{\centering \includegraphics[width=0.95\linewidth]{03_objetivo4_factor_sensorial-06-_files/figure-latex/mosaico-s3-1} 

}

\caption{Proporción de Recomendar por nivel de S3 (Apariencia).}\label{fig:mosaico-s3}
\end{figure}

\subsubsection{Dispersión con jitter: Ítems del servicio
vs.~Satisfacción}\label{dispersiuxf3n-con-jitter-uxedtems-del-servicio-vs.-satisfacciuxf3n}

\begin{figure}[H]

{\centering \includegraphics[width=0.95\linewidth]{03_objetivo4_factor_sensorial-06-_files/figure-latex/jitter-satisf-serv-1} 

}

\caption{Relación entre cada ítem del Servicio y la Satisfacción (I1).}\label{fig:jitter-satisf-serv}
\end{figure}

\subsubsection{Dispersión con jitter: Ítems del servicio
vs.~Revisitar}\label{dispersiuxf3n-con-jitter-uxedtems-del-servicio-vs.-revisitar}

\begin{figure}[H]

{\centering \includegraphics[width=0.95\linewidth]{03_objetivo4_factor_sensorial-06-_files/figure-latex/jitter-revis-serv-1} 

}

\caption{Relación entre cada ítem del Servicio y la Intención de revisitar (I2).}\label{fig:jitter-revis-serv}
\end{figure}

\subsubsection{Resumen (Servicio)}\label{resumen-servicio}

El ítem de servicio más influyente en la toma de decisiones del comensal
es \textbf{S1 (Atención)}, con \(|\rho_s|\) promedio = 0.3.

\newpage

\section{Análisis por ítems de Entorno: E1--E8 vs.~variables
dependientes}\label{anuxe1lisis-por-uxedtems-de-entorno-e1e8-vs.-variables-dependientes}

\subsection{Correlaciones de
Spearman}\label{correlaciones-de-spearman-2}

\begingroup\fontsize{9}{11}\selectfont

\begin{longtable}[t]{llcc>{}c}
\toprule
Ítem Entorno & Variable dependiente & \$\textbackslash{}rho\_s\$ & \$p\$ & Significativo\\
\midrule
\addlinespace[0.3em]
\multicolumn{5}{l}{\textbf{E1 (Decoración)}}\\
\hspace{1em}E1 (Decoración) & Satisfacción & 0.400 & < .001 & \textcolor{teal}{Sí}\\
\hspace{1em}E1 (Decoración) & Revisitar & 0.440 & < .001 & \textcolor{teal}{Sí}\\
\hspace{1em}E1 (Decoración) & Recomendar & 0.329 & < .001 & \textcolor{teal}{Sí}\\
\addlinespace[0.3em]
\multicolumn{5}{l}{\textbf{E2 (Colores)}}\\
\hspace{1em}E2 (Colores) & Satisfacción & 0.306 & < .001 & \textcolor{teal}{Sí}\\
\hspace{1em}E2 (Colores) & Revisitar & 0.368 & < .001 & \textcolor{teal}{Sí}\\
\hspace{1em}E2 (Colores) & Recomendar & 0.310 & < .001 & \textcolor{teal}{Sí}\\
\addlinespace[0.3em]
\multicolumn{5}{l}{\textbf{E3 (Mobiliario)}}\\
\hspace{1em}E3 (Mobiliario) & Satisfacción & 0.252 & < .001 & \textcolor{teal}{Sí}\\
\hspace{1em}E3 (Mobiliario) & Revisitar & 0.372 & < .001 & \textcolor{teal}{Sí}\\
\hspace{1em}E3 (Mobiliario) & Recomendar & 0.263 & < .001 & \textcolor{teal}{Sí}\\
\addlinespace[0.3em]
\multicolumn{5}{l}{\textbf{E4 (Menaje)}}\\
\hspace{1em}E4 (Menaje) & Satisfacción & 0.232 & 0.0020 & \textcolor{teal}{Sí}\\
\hspace{1em}E4 (Menaje) & Revisitar & 0.292 & < .001 & \textcolor{teal}{Sí}\\
\hspace{1em}E4 (Menaje) & Recomendar & 0.190 & 0.0116 & \textcolor{teal}{Sí}\\
\addlinespace[0.3em]
\multicolumn{5}{l}{\textbf{E5 (Manteles)}}\\
\hspace{1em}E5 (Manteles) & Satisfacción & 0.161 & 0.0326 & \textcolor{teal}{Sí}\\
\hspace{1em}E5 (Manteles) & Revisitar & 0.244 & 0.0011 & \textcolor{teal}{Sí}\\
\hspace{1em}E5 (Manteles) & Recomendar & 0.177 & 0.0190 & \textcolor{teal}{Sí}\\
\addlinespace[0.3em]
\multicolumn{5}{l}{\textbf{E6 (Espacio)}}\\
\hspace{1em}E6 (Espacio) & Satisfacción & 0.333 & < .001 & \textcolor{teal}{Sí}\\
\hspace{1em}E6 (Espacio) & Revisitar & 0.376 & < .001 & \textcolor{teal}{Sí}\\
\hspace{1em}E6 (Espacio) & Recomendar & 0.242 & 0.0012 & \textcolor{teal}{Sí}\\
\addlinespace[0.3em]
\multicolumn{5}{l}{\textbf{E7 (Asientos)}}\\
\hspace{1em}E7 (Asientos) & Satisfacción & 0.273 & < .001 & \textcolor{teal}{Sí}\\
\hspace{1em}E7 (Asientos) & Revisitar & 0.274 & < .001 & \textcolor{teal}{Sí}\\
\hspace{1em}E7 (Asientos) & Recomendar & 0.199 & 0.0080 & \textcolor{teal}{Sí}\\
\addlinespace[0.3em]
\multicolumn{5}{l}{\textbf{E8 (Música)}}\\
\hspace{1em}E8 (Música) & Satisfacción & 0.364 & < .001 & \textcolor{teal}{Sí}\\
\hspace{1em}E8 (Música) & Revisitar & 0.374 & < .001 & \textcolor{teal}{Sí}\\
\hspace{1em}E8 (Música) & Recomendar & 0.284 & < .001 & \textcolor{teal}{Sí}\\
\bottomrule
\end{longtable}
\endgroup{}

\subsubsection{Ranking de influencia
(Entorno)}\label{ranking-de-influencia-entorno}

\begin{longtable}[t]{clc}
\toprule
Ranking & Ítem Entorno & \$\&\#124;\textbackslash{}rho\_s\&\#124;\$ promedio\\
\midrule
1 & E1 (Decoración) & 0.390\\
2 & E8 (Música) & 0.341\\
3 & E2 (Colores) & 0.328\\
4 & E6 (Espacio) & 0.317\\
5 & E3 (Mobiliario) & 0.296\\
\addlinespace
6 & E7 (Asientos) & 0.249\\
7 & E4 (Menaje) & 0.238\\
8 & E5 (Manteles) & 0.194\\
\bottomrule
\end{longtable}

\newpage

\subsection{Tablas de contingencia
(Entorno)}\label{tablas-de-contingencia-entorno}

\subsubsection{E1 (Decoración) vs.~Variables
dependientes}\label{e1-decoraciuxf3n-vs.-variables-dependientes}

\begin{longtable}[t]{lrrrrrr}
\caption{\label{tab:cont-e1-satisf}E1 (Decoración) vs. Satisfacción (I1)}\\
\toprule
 & 1 & 2 & 3 & 4 & 5 & Sum\\
\midrule
1 & 2 & 4 & 5 & 0 & 0 & 11\\
2 & 0 & 7 & 4 & 5 & 0 & 16\\
3 & 0 & 10 & 4 & 5 & 1 & 20\\
4 & 2 & 10 & 19 & 23 & 0 & 54\\
5 & 2 & 8 & 15 & 41 & 9 & 75\\
\addlinespace
Sum & 6 & 39 & 47 & 74 & 10 & 176\\
\bottomrule
\end{longtable}

\begin{longtable}[t]{lrrrrrr}
\caption{\label{tab:cont-e1-revis}E1 (Decoración) vs. Revisitar (I2)}\\
\toprule
 & 1 & 2 & 3 & 4 & 5 & Sum\\
\midrule
1 & 1 & 7 & 3 & 0 & 0 & 11\\
2 & 2 & 9 & 3 & 2 & 0 & 16\\
3 & 1 & 8 & 6 & 5 & 0 & 20\\
4 & 3 & 11 & 26 & 12 & 2 & 54\\
5 & 2 & 10 & 22 & 28 & 13 & 75\\
\addlinespace
Sum & 9 & 45 & 60 & 47 & 15 & 176\\
\bottomrule
\end{longtable}

\begin{longtable}[t]{lrrr}
\caption{\label{tab:cont-e1-recom}E1 (Decoración) vs. Recomendar (I3)}\\
\toprule
 & No recomendaría & Sí recomendaría & Sum\\
\midrule
1 & 9 & 2 & 11\\
2 & 12 & 4 & 16\\
3 & 12 & 8 & 20\\
4 & 30 & 24 & 54\\
5 & 24 & 51 & 75\\
\addlinespace
Sum & 87 & 89 & 176\\
\bottomrule
\end{longtable}

\subsubsection{E2 (Colores) vs.~Variables
dependientes}\label{e2-colores-vs.-variables-dependientes}

\begin{longtable}[t]{lrrrrrr}
\caption{\label{tab:cont-e2-satisf}E2 (Colores) vs. Satisfacción (I1)}\\
\toprule
 & 1 & 2 & 3 & 4 & 5 & Sum\\
\midrule
1 & 1 & 4 & 7 & 1 & 0 & 13\\
2 & 1 & 6 & 2 & 4 & 0 & 13\\
3 & 0 & 7 & 4 & 7 & 1 & 19\\
4 & 2 & 14 & 23 & 28 & 3 & 70\\
5 & 2 & 8 & 11 & 34 & 6 & 61\\
\addlinespace
Sum & 6 & 39 & 47 & 74 & 10 & 176\\
\bottomrule
\end{longtable}

\begin{longtable}[t]{lrrrrrr}
\caption{\label{tab:cont-e2-revis}E2 (Colores) vs. Revisitar (I2)}\\
\toprule
 & 1 & 2 & 3 & 4 & 5 & Sum\\
\midrule
1 & 2 & 7 & 4 & 0 & 0 & 13\\
2 & 1 & 7 & 3 & 2 & 0 & 13\\
3 & 0 & 8 & 7 & 3 & 1 & 19\\
4 & 3 & 17 & 25 & 20 & 5 & 70\\
5 & 3 & 6 & 21 & 22 & 9 & 61\\
\addlinespace
Sum & 9 & 45 & 60 & 47 & 15 & 176\\
\bottomrule
\end{longtable}

\begin{longtable}[t]{lrrr}
\caption{\label{tab:cont-e2-recom}E2 (Colores) vs. Recomendar (I3)}\\
\toprule
 & No recomendaría & Sí recomendaría & Sum\\
\midrule
1 & 11 & 2 & 13\\
2 & 10 & 3 & 13\\
3 & 10 & 9 & 19\\
4 & 37 & 33 & 70\\
5 & 19 & 42 & 61\\
\addlinespace
Sum & 87 & 89 & 176\\
\bottomrule
\end{longtable}

\subsubsection{E3 (Mobiliario) vs.~Variables
dependientes}\label{e3-mobiliario-vs.-variables-dependientes}

\begin{longtable}[t]{lrrrrrr}
\caption{\label{tab:cont-e3-satisf}E3 (Mobiliario) vs. Satisfacción (I1)}\\
\toprule
 & 1 & 2 & 3 & 4 & 5 & Sum\\
\midrule
1 & 0 & 2 & 2 & 0 & 0 & 4\\
2 & 2 & 10 & 8 & 8 & 0 & 28\\
3 & 2 & 12 & 15 & 23 & 1 & 53\\
4 & 2 & 9 & 14 & 28 & 6 & 59\\
5 & 0 & 6 & 8 & 15 & 3 & 32\\
\addlinespace
Sum & 6 & 39 & 47 & 74 & 10 & 176\\
\bottomrule
\end{longtable}

\begin{longtable}[t]{lrrrrrr}
\caption{\label{tab:cont-e3-revis}E3 (Mobiliario) vs. Revisitar (I2)}\\
\toprule
 & 1 & 2 & 3 & 4 & 5 & Sum\\
\midrule
1 & 0 & 4 & 0 & 0 & 0 & 4\\
2 & 2 & 11 & 12 & 3 & 0 & 28\\
3 & 3 & 16 & 21 & 12 & 1 & 53\\
4 & 3 & 8 & 20 & 23 & 5 & 59\\
5 & 1 & 6 & 7 & 9 & 9 & 32\\
\addlinespace
Sum & 9 & 45 & 60 & 47 & 15 & 176\\
\bottomrule
\end{longtable}

\begin{longtable}[t]{lrrr}
\caption{\label{tab:cont-e3-recom}E3 (Mobiliario) vs. Recomendar (I3)}\\
\toprule
 & No recomendaría & Sí recomendaría & Sum\\
\midrule
1 & 3 & 1 & 4\\
2 & 20 & 8 & 28\\
3 & 29 & 24 & 53\\
4 & 25 & 34 & 59\\
5 & 10 & 22 & 32\\
\addlinespace
Sum & 87 & 89 & 176\\
\bottomrule
\end{longtable}

\subsubsection{E4 (Menaje) vs.~Variables
dependientes}\label{e4-menaje-vs.-variables-dependientes}

\begin{longtable}[t]{lrrrrrr}
\caption{\label{tab:cont-e4-satisf}E4 (Menaje) vs. Satisfacción (I1)}\\
\toprule
 & 1 & 2 & 3 & 4 & 5 & Sum\\
\midrule
1 & 0 & 1 & 1 & 0 & 0 & 2\\
2 & 1 & 7 & 5 & 9 & 0 & 22\\
3 & 1 & 13 & 13 & 14 & 1 & 42\\
4 & 3 & 13 & 18 & 32 & 3 & 69\\
5 & 1 & 5 & 10 & 19 & 6 & 41\\
\addlinespace
Sum & 6 & 39 & 47 & 74 & 10 & 176\\
\bottomrule
\end{longtable}

\begin{longtable}[t]{lrrrrrr}
\caption{\label{tab:cont-e4-revis}E4 (Menaje) vs. Revisitar (I2)}\\
\toprule
 & 1 & 2 & 3 & 4 & 5 & Sum\\
\midrule
1 & 0 & 2 & 0 & 0 & 0 & 2\\
2 & 3 & 7 & 7 & 4 & 1 & 22\\
3 & 1 & 14 & 17 & 10 & 0 & 42\\
4 & 4 & 17 & 22 & 22 & 4 & 69\\
5 & 1 & 5 & 14 & 11 & 10 & 41\\
\addlinespace
Sum & 9 & 45 & 60 & 47 & 15 & 176\\
\bottomrule
\end{longtable}

\begin{longtable}[t]{lrrr}
\caption{\label{tab:cont-e4-recom}E4 (Menaje) vs. Recomendar (I3)}\\
\toprule
 & No recomendaría & Sí recomendaría & Sum\\
\midrule
1 & 2 & 0 & 2\\
2 & 13 & 9 & 22\\
3 & 23 & 19 & 42\\
4 & 36 & 33 & 69\\
5 & 13 & 28 & 41\\
\addlinespace
Sum & 87 & 89 & 176\\
\bottomrule
\end{longtable}

\subsubsection{E5 (Manteles) vs.~Variables
dependientes}\label{e5-manteles-vs.-variables-dependientes}

\begin{longtable}[t]{lrrrrrr}
\caption{\label{tab:cont-e5-satisf}E5 (Manteles) vs. Satisfacción (I1)}\\
\toprule
 & 1 & 2 & 3 & 4 & 5 & Sum\\
\midrule
1 & 0 & 1 & 1 & 0 & 0 & 2\\
2 & 0 & 4 & 3 & 5 & 0 & 12\\
3 & 1 & 12 & 6 & 10 & 1 & 30\\
4 & 4 & 12 & 20 & 38 & 2 & 76\\
5 & 1 & 10 & 17 & 21 & 7 & 56\\
\addlinespace
Sum & 6 & 39 & 47 & 74 & 10 & 176\\
\bottomrule
\end{longtable}

\begin{longtable}[t]{lrrrrrr}
\caption{\label{tab:cont-e5-revis}E5 (Manteles) vs. Revisitar (I2)}\\
\toprule
 & 1 & 2 & 3 & 4 & 5 & Sum\\
\midrule
1 & 0 & 2 & 0 & 0 & 0 & 2\\
2 & 1 & 3 & 5 & 3 & 0 & 12\\
3 & 2 & 11 & 10 & 6 & 1 & 30\\
4 & 5 & 17 & 30 & 20 & 4 & 76\\
5 & 1 & 12 & 15 & 18 & 10 & 56\\
\addlinespace
Sum & 9 & 45 & 60 & 47 & 15 & 176\\
\bottomrule
\end{longtable}

\begin{longtable}[t]{lrrr}
\caption{\label{tab:cont-e5-recom}E5 (Manteles) vs. Recomendar (I3)}\\
\toprule
 & No recomendaría & Sí recomendaría & Sum\\
\midrule
1 & 1 & 1 & 2\\
2 & 7 & 5 & 12\\
3 & 19 & 11 & 30\\
4 & 39 & 37 & 76\\
5 & 21 & 35 & 56\\
\addlinespace
Sum & 87 & 89 & 176\\
\bottomrule
\end{longtable}

\subsubsection{E6 (Espacio) vs.~Variables
dependientes}\label{e6-espacio-vs.-variables-dependientes}

\begin{longtable}[t]{lrrrrrr}
\caption{\label{tab:cont-e6-satisf}E6 (Espacio) vs. Satisfacción (I1)}\\
\toprule
 & 1 & 2 & 3 & 4 & 5 & Sum\\
\midrule
1 & 2 & 13 & 4 & 6 & 0 & 25\\
2 & 0 & 3 & 8 & 2 & 0 & 13\\
3 & 0 & 6 & 4 & 8 & 0 & 18\\
4 & 2 & 7 & 13 & 19 & 2 & 43\\
5 & 2 & 10 & 18 & 39 & 8 & 77\\
\addlinespace
Sum & 6 & 39 & 47 & 74 & 10 & 176\\
\bottomrule
\end{longtable}

\begin{longtable}[t]{lrrrrrr}
\caption{\label{tab:cont-e6-revis}E6 (Espacio) vs. Revisitar (I2)}\\
\toprule
 & 1 & 2 & 3 & 4 & 5 & Sum\\
\midrule
1 & 3 & 15 & 5 & 2 & 0 & 25\\
2 & 0 & 5 & 7 & 1 & 0 & 13\\
3 & 2 & 4 & 7 & 5 & 0 & 18\\
4 & 2 & 8 & 16 & 14 & 3 & 43\\
5 & 2 & 13 & 25 & 25 & 12 & 77\\
\addlinespace
Sum & 9 & 45 & 60 & 47 & 15 & 176\\
\bottomrule
\end{longtable}

\begin{longtable}[t]{lrrr}
\caption{\label{tab:cont-e6-recom}E6 (Espacio) vs. Recomendar (I3)}\\
\toprule
 & No recomendaría & Sí recomendaría & Sum\\
\midrule
1 & 20 & 5 & 25\\
2 & 8 & 5 & 13\\
3 & 10 & 8 & 18\\
4 & 18 & 25 & 43\\
5 & 31 & 46 & 77\\
\addlinespace
Sum & 87 & 89 & 176\\
\bottomrule
\end{longtable}

\subsubsection{E7 (Asientos) vs.~Variables
dependientes}\label{e7-asientos-vs.-variables-dependientes}

\begin{longtable}[t]{lrrrrrr}
\caption{\label{tab:cont-e7-satisf}E7 (Asientos) vs. Satisfacción (I1)}\\
\toprule
 & 1 & 2 & 3 & 4 & 5 & Sum\\
\midrule
1 & 1 & 13 & 7 & 5 & 0 & 26\\
2 & 1 & 3 & 6 & 8 & 0 & 18\\
3 & 0 & 5 & 4 & 8 & 1 & 18\\
4 & 3 & 8 & 11 & 21 & 1 & 44\\
5 & 1 & 10 & 19 & 32 & 8 & 70\\
\addlinespace
Sum & 6 & 39 & 47 & 74 & 10 & 176\\
\bottomrule
\end{longtable}

\begin{longtable}[t]{lrrrrrr}
\caption{\label{tab:cont-e7-revis}E7 (Asientos) vs. Revisitar (I2)}\\
\toprule
 & 1 & 2 & 3 & 4 & 5 & Sum\\
\midrule
1 & 3 & 14 & 6 & 3 & 0 & 26\\
2 & 0 & 6 & 7 & 5 & 0 & 18\\
3 & 1 & 4 & 7 & 5 & 1 & 18\\
4 & 3 & 9 & 13 & 14 & 5 & 44\\
5 & 2 & 12 & 27 & 20 & 9 & 70\\
\addlinespace
Sum & 9 & 45 & 60 & 47 & 15 & 176\\
\bottomrule
\end{longtable}

\begin{longtable}[t]{lrrr}
\caption{\label{tab:cont-e7-recom}E7 (Asientos) vs. Recomendar (I3)}\\
\toprule
 & No recomendaría & Sí recomendaría & Sum\\
\midrule
1 & 20 & 6 & 26\\
2 & 10 & 8 & 18\\
3 & 8 & 10 & 18\\
4 & 20 & 24 & 44\\
5 & 29 & 41 & 70\\
\addlinespace
Sum & 87 & 89 & 176\\
\bottomrule
\end{longtable}

\subsubsection{E8 (Música) vs.~Variables
dependientes}\label{e8-muxfasica-vs.-variables-dependientes}

\begin{longtable}[t]{lrrrrrr}
\caption{\label{tab:cont-e8-satisf}E8 (Música) vs. Satisfacción (I1)}\\
\toprule
 & 1 & 2 & 3 & 4 & 5 & Sum\\
\midrule
1 & 2 & 12 & 16 & 6 & 0 & 36\\
2 & 1 & 6 & 2 & 6 & 0 & 15\\
3 & 0 & 5 & 9 & 11 & 2 & 27\\
4 & 3 & 10 & 9 & 20 & 1 & 43\\
5 & 0 & 6 & 11 & 31 & 7 & 55\\
\addlinespace
Sum & 6 & 39 & 47 & 74 & 10 & 176\\
\bottomrule
\end{longtable}

\begin{longtable}[t]{lrrrrrr}
\caption{\label{tab:cont-e8-revis}E8 (Música) vs. Revisitar (I2)}\\
\toprule
 & 1 & 2 & 3 & 4 & 5 & Sum\\
\midrule
1 & 3 & 18 & 11 & 4 & 0 & 36\\
2 & 2 & 4 & 4 & 5 & 0 & 15\\
3 & 0 & 7 & 11 & 8 & 1 & 27\\
4 & 4 & 9 & 13 & 15 & 2 & 43\\
5 & 0 & 7 & 21 & 15 & 12 & 55\\
\addlinespace
Sum & 9 & 45 & 60 & 47 & 15 & 176\\
\bottomrule
\end{longtable}

\begin{longtable}[t]{lrrr}
\caption{\label{tab:cont-e8-recom}E8 (Música) vs. Recomendar (I3)}\\
\toprule
 & No recomendaría & Sí recomendaría & Sum\\
\midrule
1 & 28 & 8 & 36\\
2 & 8 & 7 & 15\\
3 & 11 & 16 & 27\\
4 & 22 & 21 & 43\\
5 & 18 & 37 & 55\\
\addlinespace
Sum & 87 & 89 & 176\\
\bottomrule
\end{longtable}

\newpage

\subsection{Visualizaciones (Entorno)}\label{visualizaciones-entorno}

\subsubsection{Heatmap de
correlaciones}\label{heatmap-de-correlaciones-2}

\begin{figure}[H]

{\centering \includegraphics[width=0.95\linewidth]{03_objetivo4_factor_sensorial-06-_files/figure-latex/heatmap-ent-1} 

}

\caption{Correlaciones de Spearman entre los ítems del Entorno y las variables dependientes.}\label{fig:heatmap-ent}
\end{figure}

\subsubsection{Proporción de Recomendar por nivel de cada ítem del
Entorno}\label{proporciuxf3n-de-recomendar-por-nivel-de-cada-uxedtem-del-entorno}

\begin{figure}[H]

{\centering \includegraphics[width=0.95\linewidth]{03_objetivo4_factor_sensorial-06-_files/figure-latex/mosaico-ent-all-1} 

}

\caption{Proporción de Recomendar por nivel de cada ítem del Entorno.}\label{fig:mosaico-ent-all}
\end{figure}

\subsubsection{Dispersión con jitter: Ítems del entorno
vs.~Satisfacción}\label{dispersiuxf3n-con-jitter-uxedtems-del-entorno-vs.-satisfacciuxf3n}

\begin{figure}[H]

{\centering \includegraphics[width=0.95\linewidth]{03_objetivo4_factor_sensorial-06-_files/figure-latex/jitter-satisf-ent-1} 

}

\caption{Relación entre cada ítem del Entorno y la Satisfacción (I1).}\label{fig:jitter-satisf-ent}
\end{figure}

\subsubsection{Dispersión con jitter: Ítems del entorno
vs.~Revisitar}\label{dispersiuxf3n-con-jitter-uxedtems-del-entorno-vs.-revisitar}

\begin{figure}[H]

{\centering \includegraphics[width=0.95\linewidth]{03_objetivo4_factor_sensorial-06-_files/figure-latex/jitter-revis-ent-1} 

}

\caption{Relación entre cada ítem del Entorno y la Intención de revisitar (I2).}\label{fig:jitter-revis-ent}
\end{figure}

\subsubsection{Resumen (Entorno)}\label{resumen-entorno}

El ítem de entorno más influyente en la toma de decisiones del comensal
es \textbf{E1 (Decoración)}, con \(|\rho_s|\) promedio = 0.39.

\newpage

\section{Resumen general del Objetivo
4}\label{resumen-general-del-objetivo-4}

\subsection{Conclusión --- Dimensión
Producto}\label{conclusiuxf3n-dimensiuxf3n-producto}

Dentro de la dimensión Producto, el ítem más influyente es \textbf{P1
(Sabor)} (\(|\rho_s|\) promedio = 0.54), mientras que el de menor
influencia es \textbf{P3 (Frescura)} (\(|\rho_s|\) promedio = 0.376).
Las tablas de contingencia y los gráficos de barras apiladas confirman
que a mayor puntaje en los ítems de producto, la proporción de
comensales dispuestos a recomendar el establecimiento aumenta de forma
consistente. Los gráficos de dispersión muestran una relación positiva
entre los tres ítems y las variables de satisfacción y revisita.

\subsection{Conclusión --- Dimensión
Servicio}\label{conclusiuxf3n-dimensiuxf3n-servicio}

Dentro de la dimensión Servicio, el ítem más influyente es \textbf{S1
(Atención)} (\(|\rho_s|\) promedio = 0.3), y el de menor influencia es
\textbf{S3 (Apariencia)} (\(|\rho_s|\) promedio = 0.205). Los patrones
observados en las tablas de contingencia y en los gráficos de mosaico
son coherentes con los resultados de las correlaciones: los niveles más
altos de percepción de servicio se asocian con mayor proporción de
recomendación y con puntajes más elevados de satisfacción y revisita.

\subsection{Conclusión --- Dimensión
Entorno}\label{conclusiuxf3n-dimensiuxf3n-entorno}

Dentro de la dimensión Entorno, el ítem más influyente es \textbf{E1
(Decoración)} (\(|\rho_s|\) promedio = 0.39), y el de menor influencia
es \textbf{E5 (Manteles)} (\(|\rho_s|\) promedio = 0.194). La mayor
cantidad de ítems en esta dimensión (8 frente a 3 de las otras) permite
una mirada más granular: el heatmap evidencia que no todos los aspectos
del entorno pesan por igual, y las barras apiladas muestran que ciertos
elementos del ambiente tienen una relación más marcada con la decisión
de recomendar.

\subsection{Conclusión general}\label{conclusiuxf3n-general}

\begin{longtable}[t]{clcc}
\caption{\label{tab:tabla-ranking-global}Ranking global de ítems por influencia sobre las variables dependientes}\\
\toprule
Ranking & Dimensión & Ítem & \$\&\#124;\textbackslash{}rho\_s\&\#124;\$ promedio\\
\midrule
\textbf{1} & \textbf{Producto} & \textbf{P1 (Sabor)} & \textbf{0.540}\\
2 & Producto & P2 (Presentación) & 0.478\\
3 & Entorno & E1 (Decoración) & 0.390\\
4 & Producto & P3 (Frescura) & 0.376\\
5 & Entorno & E8 (Música) & 0.341\\
\addlinespace
6 & Entorno & E2 (Colores) & 0.328\\
7 & Entorno & E6 (Espacio) & 0.317\\
8 & Servicio & S1 (Atención) & 0.300\\
9 & Entorno & E3 (Mobiliario) & 0.296\\
10 & Entorno & E7 (Asientos) & 0.249\\
\addlinespace
11 & Entorno & E4 (Menaje) & 0.238\\
12 & Servicio & S2 (Confianza) & 0.228\\
13 & Servicio & S3 (Apariencia) & 0.205\\
14 & Entorno & E5 (Manteles) & 0.194\\
\bottomrule
\end{longtable}

Considerando las tres dimensiones de forma conjunta, el factor sensorial
más influyente en la toma de decisiones del comensal es \textbf{P1
(Sabor)} (dimensión \textbf{Producto}), con un \(|\rho_s|\) promedio de
0.54. Este resultado indica que, entre todos los atributos evaluados,
dicho factor presenta la asociación más fuerte y consistente con las
tres variables de resultado: satisfacción, intención de revisitar e
intención de recomendar.

La tabla de ranking global permite además identificar cuáles ítems de
cada dimensión contribuyen más a la experiencia del comensal,
información útil para priorizar esfuerzos de mejora en un
establecimiento de alimentos y bebidas.

\end{document}
